\providecommand\ExamFontPath{../../../字体/}
\documentclass[fontset=exam]{exam-zh}
\usepackage{longtable,booktabs,array,calc,multirow,fvextra}
\usepackage[normalem]{ulem}
\providecommand{\hl}[1]{#1}
\providecommand{\ul}[1]{\uline{#1}}
\DefineVerbatimEnvironment{verbatim}{Verbatim}{breaklines=true,fontsize=\small}
\providecommand{\tightlist}{\setlength{\itemsep}{0pt}\setlength{\parskip}{0pt}}
\newcommand{\pandocbounded}[1]{#1}
\setlength{\emergencystretch}{3em}
\setlength{\tabcolsep}{2pt}
\examsetup{page/size=a4paper}
\title{期末考试 A 卷}
\subject[4em]{离散数学}
\begin{document}
\maketitle
\fancyfoot[C]{第 \thepage 页（共 \pageref*{LastPage} 页）}
\pdfbookmark[1]{一、证明题（10分）}{section1}

\par\medskip{\heiti

一、证明题（10 分）

\par}\smallskip

1)(\(\neg\)P\(\land\)(\(\neg\)Q\(\land\)R))\(\lor\)(Q\(\land\)R)\(\lor\)(P\(\land\)R)\(\Leftrightarrow\)R

2)\(\exists\)x(A(x)\(\to\)B(x))\(\Leftrightarrow\) \(\forall\)xA(x)\(\to\)\(\exists\)xB(x)

\pdfbookmark[1]{二、求命题公式(P∨(Q∧R))→(P∧Q∧R)的主析取范式和主合取范式（10分）}{section4}

二、求命题公式(P\(\lor\)(Q\(\land\)R))\(\to\)(P\(\land\)Q\(\land\)R)的主析取范式和主合取范式（10 分）

\pdfbookmark[1]{三、推理证明题（10分）}{section5}

三、推理证明题（10 分）

\begin{enumerate}
\def\labelenumi{\arabic{enumi}.}
\item
\item
  C\(\lor\)D, (C\(\lor\)D)\(\to\) \(\neg\)E, \(\neg\)E\(\to\)(A\(\land\)\(\neg\)B), (A\(\land\)\(\neg\)B)\(\to\)(R\(\lor\)S)\(\Rightarrow\)R\(\lor\)S
\end{enumerate}

2) \(\forall\)x(P(x)\(\to\)Q(y)\(\land\)R(x))，\(\exists\)xP(x)\(\Rightarrow\)Q(y)\(\land\)\(\exists\)x(P(x)\(\land\)R(x))

\pdfbookmark[1]{四、设m是一个取定的正整数，证明：在任取m＋1个整数中，至少有两个整数，它们的差是m的整数倍}{section8}

四、设\emph{m}是一个取定的正整数，证明：在任取\emph{m}＋1 个整数中，至少有两个整数，它们的差是\emph{m}的整数倍

\pdfbookmark[1]{五、已知A、B、C是三个集合，证明A-(B∪C)=(A-B)∩(A-C) （15分）}{section9}

五、已知 A、B、C 是三个集合，证明 A-(B\(\cup\)C)=(A-B)\(\cap\)(A-C) （15 分）

\pdfbookmark[1]{六、已知R、S是N上的关系，其定义如下：R=<x,y>| x,y∈N∧y=x2，S=<x,y>| x,y∈N∧y=x+1}{section10}

六、已知 R、S 是 N 上的关系，其定义如下：R=\{\textless x,y\textgreater\textbar{} x,y\(\in\)N\(\land\)y=x\textsuperscript{2}\}，S=\{\textless x,y\textgreater\textbar{} x,y\(\in\)N\(\land\)y=x+1\}。求 R\textsuperscript{-1}、R*S、S*R、R\includegraphics[height=1.5ex]{figures/image7.png}\{1,2\}、S{[}\{1,2\}{]}（10 分）

\pdfbookmark[1]{七、若f:A→B和g:B→C是双射，则（gf）-1=f-1g-1（10分）。}{section11}

七、若 f:A\(\to\)B 和 g:B\(\to\)C 是双射，则（gf）\textsuperscript{-1}=f\textsuperscript{-1}g\textsuperscript{-1}（10 分）。

\pdfbookmark[1]{八、（15分）设<A，*>是半群（结合律），对A中任意元a和b，如a≠b必有a*b≠b*a，证明：}{section12}

八、（15 分）设\textless{}\emph{A}，*\textgreater 是半群（结合律），对\emph{A}中任意元\emph{a}和\emph{b}，如\emph{a}\(\ne\)\emph{b}必有\emph{a}*\emph{b}\(\ne\)\emph{b}*\emph{a}，证明：

(1)对\emph{A}中每个元\emph{a}，有\emph{a}*\emph{a}＝\emph{a}。

(2)对\emph{A}中任意元\emph{a}和\emph{b}，有\emph{a}*\emph{b}*\emph{a}＝\emph{a}。

(3)对\emph{A}中任意元\emph{a、b}和\emph{c}，有\emph{a}*\emph{b}*\emph{c}＝\emph{a}*\emph{c}。

\pdfbookmark[1]{九、给定简单无向图G＝<V，E>，且|V|＝m，|E|＝n。试证：若n≥C_m - 1^2＋2，则G是哈密尔顿图}{section16}

九、给定简单无向图\emph{G}＝\textless{}\emph{V}，\emph{E}\textgreater，且\textbar{}\emph{V}\textbar＝\emph{m}，\textbar{}\emph{E}\textbar＝\emph{n}。试证：若\emph{n}\(\ge\)\(C_{m - 1}^{2}\)＋2，则\emph{G}是哈密尔顿图

\end{document}
